\documentclass[12pt,a4paper]{article}
\usepackage[utf8]{inputenc}
\usepackage[serbian]{babel}
\usepackage{amsmath}
\usepackage{graphicx}
\usepackage{booktabs}
\usepackage{hyperref}
\usepackage{geometry}
\geometry{top=25mm, bottom=25mm, left=30mm, right=25mm}

\hypersetup{
    colorlinks=true,
    linkcolor=black,
    filecolor=black,
    urlcolor=blue,
    citecolor=black,
}

\title{\bf Pregled tema}
\author{
    {\bf Student:} Aleksa Vukadinović \quad {\bf Indeks:} 103/2022 \\[0.3cm]
    {\bf Profesor:} dr Sana Stojanović Đurđević \\[0.3cm]
    {\bf Seminarska tema:} \emph{Dark Patterns u dizajnu korisničkog interfejsa}
}
\date{}

\begin{document}

\maketitle
\newpage

\tableofcontents
\newpage

% -----------------------------------------------------------------------
\section{Tema 1 -- Istorijski pregled razvoja računarstva}

Razvoj računarstva od mehaničkih računskih mašina, preko prvih elektronskih računara,
do interneta i mobilnih platformi. Cilj teme je da se savremeni problemi sagledaju u
istorijskom kontekstu i da se razume kako su društvene okolnosti oblikovale tehnologiju
(i obrnuto).

\textbf{Predlog teme za seminarski:}
\emph{Od ENIAC-a do veštačke inteligencije: kako je svaka generacija računara menjala
društvo}

\textbf{Ključne tačke i primeri:}
\begin{itemize}
    \item Ada Lavlejs (Ada Lovelace) i prvi algoritam -- ideja programa pre postojanja
          računara.
    \item Alan Tjuring (Alan Turing) i pojam izračunljivosti; Tjuringov test kao rana
          veza računarstva i filozofije uma.
    \item Murov zakon (Moore's Law) -- eksponencijalni rast i njegove društvene posledice.
    \item Prelazak sa Web 1.0 na Web 2.0: korisnik postaje proizvođač sadržaja.
\end{itemize}

\textbf{Izvor inspiracije:}\\
\href{https://view.genial.ly/62213fa1797ddc00194500dc/interactive-content-chronological-line}%
{Interaktivni istorijski pregled razvoja računarstva (genially)}

\textbf{Primer seminarskog rada iz repozitorijuma:}\\
\href{https://github.com/sanastojanovic/RID/blob/main/rid_2023_seminarski/mi22448_Jana_Milutinovic/racunarstvo_21_vek.pdf.pdf}%
{Jana Milutinović -- \emph{Razvoj računarstva u 21. veku} (2023)}

% -----------------------------------------------------------------------
\section{Tema 2 -- Etika u računarstvu}

Etička pitanja sa kojima se susreću programeri, inženjeri i kompanije: profesionalna
odgovornost, etički kodeksi i posledice tehnoloških odluka na pojedinca i društvo.

\textbf{Predlog teme za seminarski:}
\emph{Etički kodeksi u softverskom inženjerstvu: da li ACM/IEEE smernice utiču na
realne odluke programera?}

\textbf{Ključne tačke i primeri:}
\begin{itemize}
    \item ACM/IEEE etički kodeks -- princip ,,izbegavaj nanošenje štete``.
    \item Studija slučaja: Volkswagen ,,Dieselgate`` -- softver koji vara testove
          emisije.
    \item Therac-25 -- softverska greška sa smrtnim posledicama kao primer odgovornosti.
    \item Dileme razvoja: algoritamska pristrasnost, automatizacija i gubitak radnih
          mesta.
\end{itemize}

\textbf{Izvor inspiracije:}\\
\href{https://view.genial.ly/62304dfe2abb7a00182329b0/interactive-content-etika}%
{Interaktivna prezentacija ,,Etika`` (genially)}\\
i preporučeni udžbenik \emph{Michael J.\ Quinn -- Ethics for the Information Age
(Pearson, 2017)}

\textbf{Primer seminarskog rada iz repozitorijuma:}\\
\href{https://github.com/sanastojanovic/RID/blob/main/rid_2022_seminarski/mi17116_Predrag_Mitic/predavanje-1/Odnos-etike-i-tehnoloskog-razvoja.pdf}%
{Predrag Mitić -- \emph{Odnos etike i tehnološkog razvoja} (2022)}

% -----------------------------------------------------------------------
\section{Tema 3 -- Upotreba interneta, mrežna komunikacija i spam}

Načini komunikacije preko mreže, evolucija e-pošte i servisa, kao i problem neželjenih
poruka (spam) i zloupotrebe komunikacionih kanala.

\textbf{Predlog teme za seminarski:}
\emph{Anatomija spama: od Nigerijskih prevara do modernog phishing-a kao usluge}

\textbf{Ključne tačke i primeri:}
\begin{itemize}
    \item Istorijat: prva spam poruka (1978, ARPANET) i nastanak termina ,,spam``.
    \item Tehničke odbrane: SPF, DKIM, DMARC i Bajesovi (Bayesian) filteri.
    \item Društveni inženjering: phishing, ,,spear phishing`` i lažni CEO mejlovi.
    \item Ekonomija spama -- zašto se isplati i kada je profitabilno slati milione
          poruka.
\end{itemize}

\textbf{Izvor inspiracije:}\\
\href{https://www.microsoft.com/en-us/research/publication/why-do-nigerian-scammers-say-they-are-from-nigeria/}%
{Herley, C. (2012). \emph{Why Do Nigerian Scammers Say They are from Nigeria?}
Microsoft Research}

\textbf{Primer seminarskog rada iz repozitorijuma:}\\
\href{https://github.com/sanastojanovic/RID/blob/main/rid_2022_seminarski/mi18096_Cankovic_Jelena/1_predavanje/spam_prezentacija.pdf}%
{Jelena Čanković -- \emph{Spam, internet servisi i komunikacija} (2022)}

% -----------------------------------------------------------------------
\section{Tema 4 -- Cenzura, sloboda govora i neprikladan sadržaj}

Granica između slobode govora i potrebe za moderacijom sadržaja; zaštita dece od
neprikladnog sadržaja; uloga platformi kao ,,arbitara`` javnog govora.

\textbf{Predlog teme za seminarski:}
\emph{Ko odlučuje šta sme da se kaže? Moderacija sadržaja na društvenim mrežama
između slobode govora i odgovornosti platforme}

\textbf{Ključne tačke i primeri:}
\begin{itemize}
    \item Section 230 (SAD) -- zašto platforme nisu odgovorne za sadržaj korisnika.
    \item Algoritamska vs.\ ljudska moderacija; problem ,,shadow banning``-a.
    \item Veliki Kineski zaštitni zid (Great Firewall) kao primer državne cenzure.
    \item Roditeljska kontrola, starosna verifikacija i zaštita maloletnika.
\end{itemize}

\textbf{Izvor inspiracije:}\\
\href{https://github.com/sanastojanovic/RID/blob/main/rid_2022_seminarski/mi18141_Marko_Bura/1_predavanje/cenzura_sloboda_govora_i_neprikladan_sadrzaj.pdf}%
{Seminarski rad ,,Cenzura, sloboda govora i neprikladan sadržaj`` (PDF, RID)}

\textbf{Primer seminarskog rada iz repozitorijuma:}\\
\href{https://github.com/sanastojanovic/RID/blob/main/rid_2022_seminarski/mi18141_Marko_Bura/1_predavanje/cenzura_sloboda_govora_i_neprikladan_sadrzaj.pdf}%
{Marko Bura -- \emph{Cenzura, sloboda govora i neprikladan sadržaj} (2022)}

% -----------------------------------------------------------------------
\section{Tema 5 -- Internet prevare i zavisnost od interneta}

Mehanizmi prevara na internetu i psihološki obrasci koji vode ka zavisnosti od
digitalnih platformi i društvenih mreža.

\textbf{Predlog teme za seminarski:}
\emph{Dizajn koji stvara zavisnost: kako aplikacije koriste psihologiju da nas zadrže
(i kako da se odbranimo)}

\textbf{Ključne tačke i primeri:}
\begin{itemize}
    \item Promenljive nagrade (variable rewards) -- ,,povlačenje ručice`` kao kod
          kockanja.
    \item Beskonačno skrolovanje i automatsko učitavanje (autoplay) sadržaja.
    \item Prevare: lažne nagradne igre, ,,romance scam``, lažne investicije u
          kriptovalute.
    \item Pokret ,,Time Well Spent`` i alati za digitalni wellbeing (Screen Time,
          Digital Wellbeing).
\end{itemize}

\textbf{Izvor inspiracije:}\\
\href{https://view.genial.ly/62431d6917ac64001885c51c/presentation-problemi-na-internetu}%
{Interaktivna prezentacija ,,Problemi na internetu -- prevare i zavisnost`` (genially)}\\
i dokumentarac \emph{The Social Dilemma (Netflix, 2020)}

\textbf{Primer seminarskog rada iz repozitorijuma:}\\
\href{https://github.com/sanastojanovic/RID/blob/main/rid_2022_seminarski/mi17017_Selakovic_Sara/predavanje_1/zavisnost_od_interneta.pdf}%
{Sara Selaković -- \emph{Zavisnost od interneta} (2022)}

% -----------------------------------------------------------------------
\section{Tema 6 -- Bezbednost interneta (hakovanje i malver)}

Računarska bezbednost, vrste napada i odbrane, malver, kao i etičko hakovanje kao
legitimna profesija.

\textbf{Predlog teme za seminarski:}
\emph{Etičko hakovanje i penetraciono testiranje: zašto ,,dobri momci`` moraju da
razmišljaju kao napadači}

\textbf{Ključne tačke i primeri:}
\begin{itemize}
    \item Vrste malvera: virusi, ransomware (npr.\ WannaCry, 2017), trojanci,
          špijunski softver.
    \item CIA trijada: poverljivost, integritet, dostupnost (confidentiality,
          integrity, availability).
    \item Penetraciono testiranje i Capture The Flag (CTF) takmičenja.
    \item Napadi nultog dana (zero-day) i tržište eksploita.
\end{itemize}

\textbf{Izvor inspiracije:}\\
\href{https://github.com/sanastojanovic/RID/blob/main/rid_2022_seminarski/mi18083_Andrija_Urosevic/predavanje_1/cybersec.pdf}%
{Seminarski rad o sajber bezbednosti ,,cybersec.pdf`` (RID)}\\
i predavanje \emph{Capture The Flag (CTF) kao uvod u računarsku bezbednost}

\textbf{Primer seminarskog rada iz repozitorijuma:}\\
\href{https://github.com/sanastojanovic/RID/blob/main/rid_2025_seminarski/mi19257_Katarina_Grbovic/seminarski/pentesting.pdf}%
{Katarina Grbović -- \emph{Penetraciono testiranje} (2025)}

% -----------------------------------------------------------------------
\section{Tema 7 -- Privatnost na internetu}

Zaštita ličnih podataka, praćenje korisnika, kolačići i pravna regulativa (GDPR).
Pitanje ravnoteže između personalizacije i privatnosti.

\textbf{Predlog teme za seminarski:}
\emph{Kapitalizam nadzora: kako su lični podaci postali najvrednija roba digitalne
ekonomije}

\textbf{Ključne tačke i primeri:}
\begin{itemize}
    \item Kolačići (cookies), praćenje preko više sajtova (cross-site tracking) i
          ,,fingerprinting``.
    \item GDPR (EU) i pravo na zaborav; uloga banera za saglasnost o kolačićima.
    \item Šošana Zubof (Shoshana Zuboff) i pojam ,,surveillance capitalism``.
    \item Skandal Cambridge Analytica -- zloupotreba podataka u političke svrhe.
\end{itemize}

\textbf{Izvor inspiracije:}\\
\href{https://view.genial.ly/6258686f17af9f0011537dc1/presentation-kolacici-koriscenje-i-regulativa}%
{Interaktivna prezentacija ,,Kolačići -- korišćenje i regulativa`` (genially)}\\
i knjiga \emph{Shoshana Zuboff -- The Age of Surveillance Capitalism (2019)}

\textbf{Primer seminarskog rada iz repozitorijuma:}\\
\href{https://github.com/sanastojanovic/RID/blob/main/rid_2025_seminarski/mi21098_Marko_Lazarevic/seminarski/Privacy_in_the_AI_era.pdf}%
{Marko Lazarević -- \emph{Privatnost u eri veštačke inteligencije} (2025)}

% -----------------------------------------------------------------------
\section{Tema 8 -- Informatička pismenost u Srbiji}

Sposobnost građana da razumeju, kritički procene i bezbedno koriste digitalne
tehnologije; digitalni jaz između različitih društvenih i starosnih grupa.

\textbf{Predlog teme za seminarski:}
\emph{Digitalni jaz u Srbiji: zašto je informatička pismenost pitanje društvene
pravde, a ne samo tehnologije}

\textbf{Ključne tačke i primeri:}
\begin{itemize}
    \item Razlike u digitalnim veštinama između mlađih i starijih korisnika.
    \item Ranjivost na prevare zbog nedostatka iskustva (npr.\ reklame maskirane u
          dugme ,,Preuzmi``).
    \item Uloga obrazovnog sistema i programa doživotnog učenja.
    \item Podaci Republičkog zavoda za statistiku o korišćenju IKT u domaćinstvima.
\end{itemize}

\textbf{Izvor inspiracije:}\\
Predavanje \emph{Informatička pismenost u Srbiji} (kurs 2022) i izveštaji
\href{https://www.stat.gov.rs/}{Republičkog zavoda za statistiku o upotrebi IKT}

\textbf{Primer seminarskog rada iz repozitorijuma:}\\
\href{https://github.com/sanastojanovic/RID/blob/main/rid_2022_seminarski/mi17096_Mutavd\%C5\%BEi\%C4\%87_\%C4\%90or\%C4\%91e/predavanje_1/Informaticka\%20pismenost\%20u\%20Srbiji.pdf}%
{Đorđe Mutavdžić -- \emph{Informatička pismenost u Srbiji} (2022)}

% -----------------------------------------------------------------------
\section{Pregled veza moje seminarske teme sa temama kursa}

Tema mog seminarskog rada -- \textbf{tamni obrasci (Dark Patterns)} -- direktno se
povezuje sa više tema kursa, što je u skladu sa zahtevom da rad ostvari očiglednu vezu
sa najmanje tri teme:

\begin{table}[h!]
\centering
\begin{tabular}{p{0.28\textwidth} p{0.62\textwidth}}
\toprule
\textbf{Tema kursa} & \textbf{Veza sa tamnim obrascima} \\
\midrule
2. Etika u računarstvu &
    Tamni obrasci su konkretan primer etičke dileme: dizajner svesno kreira interfejs
    protiv interesa korisnika (ACM kodeks). \\[6pt]
5. Internet prevare i zavisnost &
    Obrasci poput veštačke hitnosti i FOMO koriste iste psihološke mehanizme kao
    prevare i dizajn koji stvara zavisnost. \\[6pt]
7. Privatnost na internetu &
    Manipulativni baneri za kolačiće (,,Prihvati sve`` istaknuto) navode korisnike da
    pristanu na prikupljanje podataka. \\[6pt]
8. Informatička pismenost &
    Sposobnost prepoznavanja tamnih obrazaca direktno zavisi od digitalne pismenosti
    -- stariji korisnici su ranjiviji. \\[6pt]
6. Bezbednost interneta &
    Obrazac \emph{Sneak into Basket} često instalira neželjeni softver (McAfee, Avast)
    tokom instalacije besplatnih aplikacija. \\
\bottomrule
\end{tabular}
\end{table}

% -----------------------------------------------------------------------
\section*{Reference}

\begin{itemize}
    \item \href{https://poincare.matf.bg.ac.rs/~sana.stojanovic.djurdjevic/rid.htm}%
          {Sajt predmeta -- Računarstvo i društvo}
    \item \href{https://github.com/sanastojanovic/RID/blob/main/rid_2025_seminarski/uputstvo_2025.txt}%
          {Uputstvo za izradu seminarskog rada (2025)}
    \item \href{https://github.com/sanastojanovic/RID}{RID repozitorijum}
    \item Michael J.\ Quinn -- \emph{Ethics for the Information Age}, Pearson, 2017.
          (preporučeni udžbenik)
    \item \href{https://technologyandsociety.org/technology-and-society-magazine/}%
          {IEEE Technology and Society Magazine}
    \item \href{https://www.technologyreview.com/}{MIT Technology Review}
\end{itemize}

\end{document}
